\begin{exercise}[Numerical self-financing strategy]
\emph{Restatement.}  Consider \(d=1\), \(r=\ln(1.06)\),
\(t_0=0,t_1=1,t_2=2\), \(S^0_{t_0}=100\), and risky prices
\[
  S^1_{t_0}=100,\qquad S^1_{t_1}=120,\qquad S^1_{t_2}=100.
\]
A self-financing portfolio has
\[
  V_{t_0}=1000,\qquad V_{t_1}=1060,\qquad V_{t_2}=1042.
\]
Compute \(S^0_{t_k}\), \(\Vtilde_{t_k}\), the holdings, and compare with a
pure risk-free investment.
\end{exercise}

\begin{solution}
The risk-free asset evolves according to \(\polyref{Eq.~(1.1)}\):
\[
  S^0_{t_k}=S^0_{t_0}e^{r(t_k-t_0)}.
\]
Since \(e^r=1.06\),
\[
  S^0_{t_0}=100,\qquad
  S^0_{t_1}=106,\qquad
  S^0_{t_2}=112.36.
\]
Thus the discounted wealths are
\[
  \Vtilde_{t_0}=\frac{1000}{100}=10,\qquad
  \Vtilde_{t_1}=\frac{1060}{106}=10,\qquad
  \Vtilde_{t_2}=\frac{1042}{112.36}\simeq 9.274.
\]
The discounted risky prices are
\[
  \Stilde^1_{t_0}=1,\qquad
  \Stilde^1_{t_1}=\frac{120}{106}\simeq 1.132,
  \qquad
  \Stilde^1_{t_2}=\frac{100}{112.36}\simeq 0.890.
\]

Because the portfolio is self-financing,
\[
  \Vtilde_{t_k}-\Vtilde_{t_{k-1}}
    =
  \theta^1_{t_k}
  \bigl(\Stilde^1_{t_k}-\Stilde^1_{t_{k-1}}\bigr),
\]
since the discounted risk-free asset is identically \(1\).  For the first
period,
\[
  10-10
    =
  \theta^1_{t_1}
  \left(\frac{120}{106}-1\right),
\]
so
\[
  \theta^1_{t_1}=0.
\]
The initial wealth identity gives
\[
  1000=\theta^0_{t_1}S^0_{t_0}+\theta^1_{t_1}S^1_{t_0}
       =100\theta^0_{t_1},
\]
therefore
\[
  \theta^0_{t_1}=10.
\]

For the second period,
\[
  \theta^1_{t_2}
  =
  \frac{\Vtilde_{t_2}-\Vtilde_{t_1}}
       {\Stilde^1_{t_2}-\Stilde^1_{t_1}}
  =
  \frac{1042/112.36-10}{100/112.36-120/106}
  =3.
\]
Using \(V_{t_1}=1060\), the risk-free holding at \(t_2\) is obtained from
\[
  1060=\theta^0_{t_2}S^0_{t_1}+\theta^1_{t_2}S^1_{t_1}
      =106\theta^0_{t_2}+3\cdot 120.
\]
Hence
\[
  \theta^0_{t_2}=\frac{700}{106}=\frac{350}{53}\simeq 6.604.
\]
The terminal check is
\[
  \theta^0_{t_2}S^0_{t_2}+\theta^1_{t_2}S^1_{t_2}
    =
  \frac{350}{53}\cdot 112.36+3\cdot 100
    =1042.
\]

A strategy investing only in the risk-free asset and starting from \(1000\)
would have value
\[
  1000(1.06)^2=1123.60
\]
at \(t_2\).  The given investor matches the risk-free account over the first
period, then holds \(3\) shares of the risky asset over the second period.
Since the risky asset falls from \(120\) to \(100\), the final value is
\[
  1042<1123.60.
\]
The strategy therefore underperforms the pure risk-free strategy by
\(81.60\).
\end{solution}
